% opencmdb — Administrator Manual. Build: see ../Makefile (LuaLaTeX; TEXINPUTS -> ../common).
\documentclass[11pt,oneside,listof=totoc,bibliography=totoc]{scrreprt}
\usepackage{opencmdb-manual}
\setmanualname{Administrator Manual}

\begin{document}
\makemanualtitle{Administrator Manual}{Installing, configuring, securing, and operating opencmdb}

\tableofcontents

\chapter{Introduction}
This manual is for the person who \emph{runs} opencmdb: installing it, configuring its data
sources, securing it, and keeping it backed up. If you are looking to use opencmdb day to
day — reviewing discoveries and documenting devices — see the companion \textbf{User Manual}.

\begin{note}
opencmdb ships as a \textbf{single binary}. There is no application server to manage, no
external job queue, and no second database to keep in sync.
\end{note}

\chapter{Architecture at a glance}
\begin{itemize}
  \item \textbf{Language / runtime:} Rust, one statically-linked binary.
  \item \textbf{Database:} MariaDB 10.11+ is the \emph{only} supported engine. SQLite and
        MySQL are not supported; PostgreSQL is not supported at this stage.
  \item \textbf{Web stack:} server-rendered HTML with lightweight interactivity; polling
        rather than server-sent events; an internal scheduler (no cron, Redis, or external
        workers).
  \item \textbf{Discovery:} a zero-privilege UniFi connector and a generic ARP/ping scanner,
        each isolated behind a contract-tested connector interface.
\end{itemize}

\begin{warning}
MariaDB is a deliberate, load-bearing choice — it is the engine included in Synology DSM's
automatic backups, so your opencmdb data is covered by the NAS backup you already run. Do
not attempt to point opencmdb at SQLite, MySQL, or PostgreSQL.
\end{warning}

\chapter{Requirements}
\begin{itemize}
  \item A host running Docker — a Synology NAS (Container Manager) is the priority target; a
        native binary is also supported.
  \item \textbf{MariaDB 10.11.11} (the exact version validated in CI, matching the DSM~7
        package, so development, CI, and production run the same engine).
  \item Outbound access from opencmdb to your discovery sources (e.g.\ the UniFi controller).
\end{itemize}

\chapter{Installation}
\begin{planned}
Step-by-step installation — the Docker image, a Synology Container Manager walkthrough, and
a native-binary path — will be documented here once the first release image is published to
Docker Hub. The outline below reflects the intended deployment.
\end{planned}

\section{Docker (Synology Container Manager)}
The intended flow: pull the published image, provide a MariaDB connection and an encryption
key, map a data volume, and start the container behind your reverse proxy.

\section{Container networking}
\label{sec:container-networking}
Leave the container in its \emph{own} network namespace --- the default for any ordinary Docker
network. This matters more than it looks.

\begin{warning}
\textbf{Do not use \texttt{network\_mode: host}.} It is the intuitive choice for a network
scanner and it is the wrong one, because it \emph{removes} a permission instead of granting one.

Docker sets \texttt{net.ipv4.ping\_group\_range} to \texttt{0 2147483647} inside a container's
own network namespace. That is precisely what allows opencmdb to open an unprivileged ICMP
datagram socket and scan as a non-root user, with no capability whatsoever. Host networking
shares the \emph{host's} namespace and inherits the host's value instead --- and many hosts,
Synology DSM among them, ship \texttt{1 0}, an empty range. The socket is then refused.

The failure is quiet: scanning is best-effort, so opencmdb logs a warning and carries on
serving. The container looks healthy, \texttt{/healthz} returns 200, and the product observes
nothing at all. Host mode additionally costs you port isolation --- a collision with an existing
service surfaces only as \texttt{Address in use} late in startup --- and leaves the container
without an address on any Docker network, which breaks reverse-proxy discovery.
\end{warning}

ICMP echo is routed and crosses a NAT without difficulty, so an ordinary bridge network scans a
LAN correctly. If you would rather reach opencmdb at a fixed address with no port mapping,
attach it to a \texttt{macvlan} or \texttt{ipvlan} network: it keeps a separate namespace, so
ICMP still works, while gaining real layer-2 presence --- which is also what the future ARP and
MAC discovery will want.

\begin{note}
With \texttt{macvlan}, a container normally cannot reach the address of the host interface it is
attached to. If your database runs on that same host, connect to a \emph{different} interface of
it, or keep an ordinary bridge network.
\end{note}

\begin{note}
Should you need host networking anyway, the published image carries the
\texttt{cap\_net\_raw} file capability on its binary, so the scan works there too --- as the
non-root user, with no host \texttt{sysctl} to set and persist --- provided the container is
granted \texttt{NET\_RAW} (\texttt{cap\_add} in the compose file, which the reference deployment
already does). Without that file capability, a non-root process loses the container's
capabilities at \texttt{exec} and the socket is refused regardless.
\end{note}

\section{Native binary}
The same single binary can run directly on a Linux host, reading its configuration from a
file and the environment.

\begin{note}
Running natively, opencmdb is subject to the host's own
\texttt{net.ipv4.ping\_group\_range}. If it is empty (\texttt{1 0}) the scan cannot open its
socket; either widen it or grant the binary the \texttt{cap\_net\_raw} file capability.
\end{note}

\chapter{Configuration}
opencmdb is configured through a layered configuration file plus environment overrides. The
database connection is given as five separate values, and opencmdb assembles the connection URL
itself:

\begin{lstlisting}[language=bash]
# Example — placeholder values only (use your own).
DATABASE_HOST     = "192.0.2.5"
DATABASE_PORT     = "3306"
DATABASE_NAME     = "opencmdb"
DATABASE_USERNAME = "opencmdb"
DATABASE_PASSWORD = "CHANGE_ME"
\end{lstlisting}

Write the password \emph{exactly as it is}. Characters that are reserved inside a URL ---
\texttt{@}, \texttt{:}, \texttt{/}, \texttt{\#}, \texttt{?}, \texttt{\%}, spaces --- need no
treatment at all here, because opencmdb percent-encodes the user info when it builds the URL.
Having to do that by hand, and the silent authentication failures that followed from forgetting
it, is exactly why the connection is no longer configured as a single string.

\begin{warning}
There is one mangling opencmdb cannot protect you from, because it happens \emph{before} the
process starts: \textbf{a \texttt{\$} in the password must be doubled --- \texttt{\$\$}}. Docker
Compose interpolates the contents of your \texttt{.env} file, so a \texttt{\$} begins what it
takes to be a variable reference and the rest of the value is discarded. A password written
\texttt{pa\$word} reaches opencmdb as \texttt{pa}, and you get an opaque ``access denied''.
Measured: \texttt{abc\$def} arrives as \texttt{abc}; \texttt{abc\$\$def} arrives as
\texttt{abc\$def}.
\end{warning}

\begin{note}
A single \texttt{DATABASE\_URL} is still honoured when none of the \texttt{DATABASE\_*}
variables is set, so existing deployments keep working, but it is \textbf{deprecated} and
opencmdb warns on startup when it is used. With it, the manual percent-encoding above is back:
\texttt{@}~$\rightarrow$~\texttt{\%40}, \texttt{:}~$\rightarrow$~\texttt{\%3A},
\texttt{/}~$\rightarrow$~\texttt{\%2F}, \texttt{\#}~$\rightarrow$~\texttt{\%23},
\texttt{?}~$\rightarrow$~\texttt{\%3F}, \texttt{\%}~$\rightarrow$~\texttt{\%25},
space~$\rightarrow$~\texttt{\%20}.
\end{note}

\begin{warning}
\texttt{DATABASE\_HOST} is the address of your database server \emph{as seen from inside the
container}. It is not \texttt{127.0.0.1}: the container has its own network namespace, so
loopback means the container itself.
\end{warning}

\begin{note}
Use RFC~5737 documentation addresses (\texttt{192.0.2.0/24}) and example hostnames in your
own notes and screenshots — never paste real network data into shared documents.
\end{note}

\begin{planned}
The full configuration reference (every key, its default, and its effect) will be generated
from the code once the configuration surface is implemented.
\end{planned}

\chapter{The database}
\section{Provisioning}
Create a dedicated database and user for opencmdb on your MariaDB 10.11 instance. On a
Synology NAS, this is the DSM-managed MariaDB package. \textbf{Nothing in the image or the
compose file does this for you} --- it is a prerequisite of the first start, and forgetting it
is the most common cause of a failed first deployment:

\begin{lstlisting}[language=sql]
CREATE DATABASE opencmdb CHARACTER SET utf8mb4 COLLATE utf8mb4_bin;
CREATE USER 'opencmdb'@'%' IDENTIFIED BY 'your-password';
GRANT ALL PRIVILEGES ON opencmdb.* TO 'opencmdb'@'%';
FLUSH PRIVILEGES;
\end{lstlisting}

The binary collation is required: identity comparison must be exact and must never depend on
the database's locale. Narrow the \texttt{'\%'} host to the address opencmdb actually connects
from once the connection works --- but read the following before guessing what that address is.

\begin{warning}
\textbf{A grant matches the address the server \emph{sees}, not the address you dial.} When
authentication fails, MariaDB names the exact identity it rejected:
\texttt{Access denied for user 'opencmdb'@'\textit{host}'}. Grant \emph{that} \textit{host}.

On a multi-homed machine the two differ, and the effect is thoroughly disorienting: traffic sent
to one interface can leave by another, so MariaDB matches on the source address it observes ---
or on the name that address reverse-resolves to --- rather than on the one you configured. The
tell-tale sign is that the host in the error message \emph{changes} as you change the URL. Grant
every address the client can present, or use \texttt{'\%'} while you are still finding out.
\end{warning}

\section{Migrations}
Schema migrations are embedded in the binary and applied on startup; you do not run them by
hand. Every text column uses a binary collation so that identity comparison is exact and
never depends on the database's locale — a correctness rule enforced in CI.

\chapter{Configuring discovery sources}
\section{UniFi}
Provide the controller URL and an API key. The connector reads devices, IPs, ports, SSIDs,
VLANs, and DHCP leases without elevated privileges.

\section{ARP/ping}
Declare the subnets to scan and a cadence within allowed bounds; you can also trigger an
on-demand scan. Where the host lacks the raw-socket capability, the scanner falls back to
ping-only.

Probes overlap: up to 64 are in flight at once by default
(\texttt{OPENCMDB\_SCAN\_CONCURRENCY}). This is a \emph{politeness} bound rather than a
throughput one --- a probe sends sixteen bytes and then waits, so the scan is single-threaded
whatever you set, and raising the limit costs no CPU. What it does raise is the load on your
gateway's ARP table, which is why a \texttt{/22} does not simply fire all 1022 probes at once.
Lower it on a fragile network.

\begin{note}
The wall-clock cost of a scan is driven by the addresses that \emph{do not} answer: each one
waits out the full probe timeout. With overlapping probes a mostly-empty \texttt{/24} completes
in seconds; scanning strictly one address at a time
(\texttt{OPENCMDB\_SCAN\_CONCURRENCY=1}) would take several minutes for the same result.
\end{note}

\begin{planned}
Source configuration screens and the per-source cadence controls are specified in the UX
design and land with the respective connectors.
\end{planned}

\chapter{Security}
\section{Threat model}
opencmdb is designed to protect against a \emph{stolen database or backup} and against
\emph{unauthenticated network access} — not against a local root attacker on the host.

\section{Secrets at rest}
Source API keys and other secrets are encrypted at rest. The encryption key is kept
\textbf{separate from the database volume}, so a stolen database file alone does not reveal
your source credentials. API-key rotation and an encrypted secret backup are supported.

\section{Transport security}
Terminate TLS at a reverse proxy in front of opencmdb; TLS termination is a deployment
concern, not something opencmdb does itself.

\begin{warning}
Never store the encryption key on the same volume as the database. A backup that contains
both defeats the entire at-rest protection.
\end{warning}

\chapter{Operations}
\section{Health and observability}
opencmdb exposes Prometheus metrics and a source-health view: for each source and scope, its
liveness (\texttt{live}/\texttt{blind} with a named cause) and any capability downgrades,
reported as notifiable events. A ``what changed since last visit'' summary highlights new and
conflicting items over routine churn.

\section{Scheduling}
Discovery runs on an internal scheduler within the single binary. There is no external cron
or job runner to configure.

\begin{planned}
The metrics catalogue and the operational dashboard are specified; the slice shipping first
covers source health and the ``what changed'' summary.
\end{planned}

\chapter{Backup and restore}
Because opencmdb stores its state in MariaDB, and MariaDB is included in DSM's automatic
backups, your primary backup is the one your NAS already performs. Keep the separate
encryption key backed up \emph{independently} of the database so a restore can decrypt
secrets.

\begin{planned}
A tested backup-and-restore runbook (including the encryption-key handling) will be provided
before the first stable release.
\end{planned}

\chapter{Upgrades}
Upgrades are performed by replacing the binary or the container image; embedded migrations
bring the schema forward on startup. Cargo.lock is committed and builds are reproducible, so
a given release always resolves the same dependencies.

\appendix
\chapter{Troubleshooting}
The source-health view is the first place to look once opencmdb is running: a scope that has
gone \texttt{blind} explains why its data is frozen and its alerts are silent.

Getting to that point is another matter. The table below collects the failures observed during
real first deployments --- none of them faults in the reconciliation engine, all of them
obstacles between a running MariaDB and a scanning opencmdb.

\begin{note}
A fatal startup error is logged in full --- its whole cause chain --- to standard output
\emph{and} to the daily log file, and the process then exits non-zero. If a container is
restarting in a loop, \texttt{opencmdb.YYYY-MM-DD.log} says why.
\end{note}

\begin{center}
\small
\begin{tabular}{@{}p{0.36\linewidth}p{0.28\linewidth}p{0.30\linewidth}@{}}
\toprule
\textbf{Symptom} & \textbf{Cause} & \textbf{Fix} \\
\midrule
\texttt{1045 Access denied for user 'opencmdb'@'\ldots'}
  & No grant for the identity MariaDB sees
  & Grant exactly the host named in the message \\
\addlinespace
\ldots\ and that host changes when you change the address
  & Multi-homed database server: replies leave by another interface
  & Grant every address the client can present \\
\addlinespace
\ldots\ and the password is definitely correct
  & A \texttt{\$} was truncated by Docker Compose
  & Double it: \texttt{\$\$} \\
\addlinespace
\texttt{Address in use (os error 98)}
  & Another service already holds the port
  & Change the host side of \texttt{ports:} \\
\addlinespace
\texttt{startup scan failed \ldots\ could not open an ICMP socket}
  & Host networking inherited an empty \texttt{ping\_group\_range}
  & Drop \texttt{network\_mode: host} (\S\ref{sec:container-networking}) \\
\addlinespace
The page loads but shows no observed data
  & The scan failed with a non-fatal warning
  & Look for \texttt{startup scan failed} in the log \\
\addlinespace
The page shows no gap at all
  & Observed state exists, but nothing is declared yet
  & Declare an entity carrying an \texttt{ipv4} \\
\bottomrule
\end{tabular}
\end{center}

\end{document}
